%%%%%%%%%%%%%%%%%%%%%%%%%%%%%%%%%%%%%%%%%
% UNSW Formal Letter
% LaTeX Template
% Version 1.0 (10/10/19)
%
% This template is based upon a LaTeX letter template by http://www.LaTeXTemplates.com 
%
% Authors:
% Brian Moses
% Vel (vel@LaTeXTemplates.com)
% Gordon Stevenson (gordon.stevenson@unsw.edu.au)
%
%
% TO DO: move up opening and UNSW Sydney Logo?
%
% License:
% CC BY-NC-SA 3.0 (http://creativecommons.org/licenses/by-nc-sa/3.0/)
%
%%%%%%%%%%%%%%%%%%%%%%%%%%%%%%%%%%%%%%%%%

%----------------------------------------------------------------------------------------
%	PACKAGES AND OTHER DOCUMENT CONFIGURATIONS
%----------------------------------------------------------------------------------------

\documentclass[10pt, a4paper]{letter} % Set the font size (10pt, 11pt and 12pt) and paper size (letterpaper, a4paper, etc)

\input{structure.tex} % Include the file that specifies the document structure

\longindentation=0pt % Un-commenting this line will push the closing "Sincerely," and date to the left of the page

\setlength{\parindent}{0cm}

\usepackage{fancyhdr}
\pagestyle{fancy}
\fancyhf{}
\renewcommand{\headrulewidth}{0pt}
\fancyfoot[CE,CO]{\footnotesize {SCHOOL OF PUBLIC HEALTH AND COMMUNITY MEDICINE, UNSW SYDNEY NSW 2052 AUSTRALIA \\
T +61 (2) 9385 2517 | F +61 (2) 9313 6185 | ABN 57 195 873 179 | CRICOS Provider Code 00098G} }



%----------------------------------------------------------------------------------------
%	YOUR INFORMATION
%----------------------------------------------------------------------------------------

\Who{Sender \\} % Your name

\Title{Sender Title \\ Personal contact number/s \\ Reference (if needed)} % Your title, leave blank for no title

%----------------------------------------------------------------------------------------
%	HEADER CONTENTS
%----------------------------------------------------------------------------------------

\logo{unsw_syd} % Logo filename, your logo should have square dimensions (i.e. roughly the same width and height), if it does not, you will need to adjust spacing within the HEADER STRUCTURE block in structure.tex (read the comments carefully!)

%----------------------------------------------------------------------------------------

\begin{document}

\begin{letter}{Addressee’s Name \\
Title \\
Company Name \\
Street or Postal Address \\
City State Post Code \\
Country}

\opening{Salutation}

This is an example of a typical letter format (set in Arial 10). The layout is consistent with modern business practice and is designed to use as few key strokes as possible (absenting all punctuation other than in text). It is recommended that characters are set flush left, with text lines ending in alignment with the centre of the address and Symbol block at the right (set a 32 mm left margin and a 40 mm right margin). Insert your name and any contact numbers at the end of your document, ranged left, underneath your signature or in the footer. Ensure you personalise the footer details with your faculty, school, centre or unit contact details. 

For most letters the date should appear first, in alignment with the bottom of the logo; usually two lines are adequate between date and addressee’s name and address, and between the address and salutation. Paragraphs always start flush left; never indent the first line; leave one line between all paragraphs. Usually two lines are left between the last line of the last paragraph and a closing further four to the sender’s name to allow space for signing. \\[2\baselineskip]

The recommended typeface for use with the UNSW letterhead is Arial. 
\\[2\baselineskip]

\closing{Complimentary close, ,}

\end{letter}

\end{document}